\documentclass[letterpaper]{article}
\usepackage[margin=1in]{geometry}
\usepackage{times}
\usepackage{helvet}
\usepackage{courier}
\usepackage{enumitem}
\usepackage{xcolor}

\setlength{\parindent}{0pt}
\setlength{\parskip}{6pt}

\newcommand{\question}[3]{%
    \textbf{#1} \textit{#2}\\
    \textcolor{blue}{#3}\par
}

\title{Reproducibility Checklist}
\author{Training Leaves Traces: Centered Residual-Product Signatures\\for White-Box Model Lineage Verification}
\date{}

\begin{document}
\maketitle

\section*{Instructions for Authors}
This document outlines key aspects for assessing reproducibility. Responses are provided below each question.

\section{General Paper Structure}

\question{1.1. Includes a conceptual outline and/or pseudocode description of AI methods introduced}
{(yes/partial/no/NA)}
{yes -- Section 3 provides the mathematical formulation of diagonal-dominance scoring (Eq. 1-3), Section 4 details the verification protocol with equations for centered signatures (Eq. 4), Hungarian matching (Eq. 5), and lineage score (Eq. 6-8). Table 2 provides architecture-aware factorization rules.}

\question{1.2. Clearly delineates statements that are opinions, hypothesis, and speculation from objective facts and results}
{(yes/no)}
{yes -- Hypotheses are clearly stated (e.g., ``gradient coupling hypothesis'' in Section 3.1), empirical results are presented separately in Section 5, and limitations are explicitly discussed in the Limitations section.}

\question{1.3. Provides well-marked pedagogical references for less-familiar readers to gain background necessary to replicate the paper}
{(yes/no)}
{yes -- References to foundational work on residual networks (He et al. 2016), CKA (Kornblith et al. 2019), SVCCA (Raghu et al. 2017), watermarking (Uchida et al. 2017, Adi et al. 2018), and model merging (Wortsman et al. 2022) are provided throughout.}

\section{Theoretical Contributions}

\question{2.1. Does this paper make theoretical contributions?}
{(yes/no)}
{yes}

\question{2.2. All assumptions and restrictions are stated clearly and formally}
{(yes/partial/no)}
{yes -- Section 2.1 formally defines white-box model lineage verification, stating requirements (compatible residual architectures, same depth $L$ and hidden dimension $d$) and what constitutes ``related'' vs ``unrelated'' checkpoints.}

\question{2.3. All novel claims are stated formally (e.g., in theorem statements)}
{(yes/partial/no)}
{yes -- Proposition 1 (correct-pair score formula, Eq. 5) and Corollary 1 (mismatched baseline) are formally stated with proofs. The null model corollary in Appendix A.1 provides formal guarantees.}

\question{2.4. Proofs of all novel claims are included}
{(yes/partial/no)}
{yes -- Proof of Proposition 1 is given inline (Section 3.2). Corollary derivations follow from Frobenius inner product orthogonality. Additional proofs in Appendix A.}

\question{2.5. Proof sketches or intuitions are given for complex and/or novel results}
{(yes/partial/no)}
{yes -- Section 3.1 provides detailed intuition for why gradient coupling creates diagonal dominance, including the geometric interpretation (cosine of angle between $M$ and identity). Figure 1 visualizes the mechanism.}

\question{2.6. Appropriate citations to theoretical tools used are given}
{(yes/partial/no)}
{yes -- Hungarian algorithm (Kuhn 1955), Frobenius norm properties, and dynamical isometry (Pennington et al. 2017) are cited.}

\question{2.7. All theoretical claims are demonstrated empirically to hold}
{(yes/partial/no/NA)}
{yes -- Section 5.1 tests all three predictions: (1) fingerprint is learned not present at init, (2) requires skip connection, (3) arises from gradient coupling not isometry. Each claim has dedicated experiments with quantitative results.}

\question{2.8. All experimental code used to eliminate or disprove claims is included}
{(yes/no/NA)}
{yes -- Code for all experiments including ablations (init\_scheme\_ablation.py, rq1\_gradient\_coupling.py) will be released. The gradient shuffling ablation (Part B) and synthetic injection (Part C) directly test necessity and sufficiency of gradient coupling.}

\section{Dataset Usage}

\question{3.1. Does this paper rely on one or more datasets?}
{(yes/no)}
{yes}

\question{3.2. A motivation is given for why the experiments are conducted on the selected datasets}
{(yes/partial/no/NA)}
{yes -- CIFAR-10 is used for controlled MLP/ResNet experiments (standard benchmark, enables fast iteration). Public HuggingFace checkpoints (GPT-2, LLaMA-2, ViT, Whisper, ResNet) are used to demonstrate real-world applicability across language/vision/speech domains.}

\question{3.3. All novel datasets introduced in this paper are included in a data appendix}
{(yes/partial/no/NA)}
{NA -- No novel datasets are introduced. The paper uses existing public datasets (CIFAR-10, ImageNet-pretrained models) and generates synthetic regression targets with documented seeds.}

\question{3.4. All novel datasets introduced in this paper will be made publicly available upon publication of the paper with a license that allows free usage for research purposes}
{(yes/partial/no/NA)}
{NA -- No novel datasets introduced.}

\question{3.5. All datasets drawn from the existing literature (potentially including authors' own previously published work) are accompanied by appropriate citations}
{(yes/no/NA)}
{yes -- CIFAR-10 (Krizhevsky 2009), HuggingFace (2024), GPT-2 (Radford et al. 2019), BERT (Devlin et al. 2019), LLaMA-2 (Touvron et al. 2023), Mistral (Jiang et al. 2023), Qwen (Yang et al. 2024), DeepSeek (2025), ViT (Dosovitskiy et al. 2021), ResNet (He et al. 2016), Whisper (Radford et al. 2023) are all cited.}

\question{3.6. All datasets drawn from the existing literature (potentially including authors' own previously published work) are publicly available}
{(yes/partial/no/NA)}
{yes -- CIFAR-10 is publicly available via torchvision. All model checkpoints are publicly available on HuggingFace (links provided in code).}

\question{3.7. All datasets that are not publicly available are described in detail, with explanation why publicly available alternatives are not scientifically satisficing}
{(yes/partial/no/NA)}
{NA -- All datasets used are publicly available.}

\section{Computational Experiments}

\question{4.1. Does this paper include computational experiments?}
{(yes/no)}
{yes}

\question{4.2. This paper states the number and range of values tried per (hyper-) parameter during development of the paper, along with the criterion used for selecting the final parameter setting}
{(yes/partial/no/NA)}
{partial -- The paper uses standard hyperparameters (Adam with lr=1e-3, gradient clipping=1.0) without extensive tuning. The lineage threshold $\tau=3.0$ is derived from statistical calibration (FPR $<$ 0.13\% under Gaussian null). Architecture parameters (depth, hidden dim) are specified for each experiment in the Reproducibility Statement and Appendix B.}

\question{4.3. Any code required for pre-processing data is included in the appendix}
{(yes/partial/no)}
{yes -- Data preprocessing is minimal (standard CIFAR-10 normalization via torchvision transforms, model loading via HuggingFace). All preprocessing code is included in the experiment scripts (e.g., get\_cifar\_loaders function).}

\question{4.4. All source code required for conducting and analyzing the experiments is included in a code appendix}
{(yes/partial/no)}
{yes -- 83 Python scripts covering all experiments are included: lineage\_phase1\_mlp.py (MLP benchmark), rq1\_gradient\_coupling.py (mechanism tests), llm\_lineage\_real.py (LLM validation), laundering\_llm.py (laundering experiments), plus baseline implementations (CKA, SVCCA, IPGuard regression analogue, aligned Frobenius, SVD distance, weight cosine).}

\question{4.5. All source code required for conducting and analyzing the experiments will be made publicly available upon publication of the paper with a license that allows free usage for research purposes}
{(yes/partial/no)}
{yes -- Code, configuration files, baseline implementations, per-pair score JSONs, and plotting scripts will be released upon acceptance under an open-source license.}

\question{4.6. All source code implementing new methods have comments detailing the implementation, with references to the paper where each step comes from}
{(yes/partial/no)}
{partial -- Core functions include docstrings explaining the computation (e.g., diagonal\_dominance function documents $s = |tr(M)| / ||M||_F$). Equation references to the paper are included in key functions. Further inline comments reference specific paper sections.}

\question{4.7. If an algorithm depends on randomness, then the method used for setting seeds is described in a way sufficient to allow replication of results}
{(yes/partial/no/NA)}
{yes -- All experiments use explicit seeds. The Reproducibility Statement specifies: ``All results in this paper are produced by deterministic scripts with fixed seeds.'' Seeds are passed as arguments (e.g., --seeds 0,1,2) and set via torch.manual\_seed() and np.random.seed(). Synthetic data generation uses torch.Generator().manual\_seed(seed).}

\question{4.8. This paper specifies the computing infrastructure used for running experiments (hardware and software), including GPU/CPU models; amount of memory; operating system; names and versions of relevant software libraries and frameworks}
{(yes/partial/no)}
{partial -- The paper mentions GPU experiments and CPU fallback (device=``cuda'' if torch.cuda.is\_available() else ``cpu''). LLM experiments specify memory constraints (16GB hosts, ml.g6.xlarge with 24GB GPU). Software dependencies: PyTorch, NumPy, SciPy, pandas, transformers (HuggingFace). Exact versions will be pinned in released code.}

\question{4.9. This paper formally describes evaluation metrics used and explains the motivation for choosing these metrics}
{(yes/partial/no)}
{yes -- Metrics are formally defined: diagonal-dominance ratio $s(M)$ (Eq. 3), lineage score $\mathcal{L}$ (Eq. 6), z-score calibration (Eq. 7-8), AUROC for binary classification, Gap-Z for distillation rejection margin. Motivation: $s(M)$ captures trace concentration (identity-component energy), $\mathcal{L}$ aggregates block-level similarity, AUROC measures discrimination.}

\question{4.10. This paper states the number of algorithm runs used to compute each reported result}
{(yes/no)}
{yes -- The Reproducibility Statement specifies: 3 seeds for MLP experiments (depth-24), 52-pair and 159-pair benchmarks with exact construction (e.g., ``2 trained reference checkpoints, each paired with 5 related-checkpoint kinds $\times$ 6 derivative seeds = 30 truly-related pairs''). Table 3 notes $n{=}75$ related, $n{=}84$ unrelated.}

\question{4.11. Analysis of experiments goes beyond single-dimensional summaries of performance (e.g., average; median) to include measures of variation, confidence, or other distributional information}
{(yes/no)}
{yes -- Results include mean and min scores (Table 3: mean $\mathcal{L}{=}0.99$, min 0.97), standard deviation for null calibration ($\sigma_0$ in Eq.\ 7), confidence intervals (FPR $<$ 0.13\% at $\tau{=}3.0$), score ranges per model family (Table 4), and distributional plots (Figure 1c, Appendix Figure 4).}

\question{4.12. The significance of any improvement or decrease in performance is judged using appropriate statistical tests (e.g., Wilcoxon signed-rank)}
{(yes/partial/no)}
{partial -- Z-score calibration against null distribution (Eq.\ 7) provides statistical grounding. AUROC with 95\% confidence threshold ($\tau{=}1.645$) and high-stakes threshold ($\tau{=}3.0$, FPR $<$ 0.13\%) are used. Gap-Z (Table 5) measures standardized separation. Spearman correlation ($\rho$) quantifies monotonic relationships in harder-regime benchmarks.}

\question{4.13. This paper lists all final (hyper-)parameters used for each model/algorithm in the paper's experiments}
{(yes/partial/no/NA)}
{yes -- Appendix B (Per-Architecture Details) specifies all parameters. MLP benchmark: depth$=$16/24, in\_dim$=$16/24, hidden$=$48/64. Training: Adam lr$=$1e-3, epochs$=$200, batch$=$256, grad\_clip$=$1.0. Fine-tuning: epochs$=$50, lr$=$3e-4. Noise: $\sigma_{\mathrm{rel}} \in [0.01, 0.15]$. Pruning: 10--85\%. Quantization: 16--256 levels. LLM experiments: SwiGLU factorization $M = W_{\mathrm{down}} W_{\mathrm{up}}$, 32 layers, $d{=}4096$.}

\end{document}
